\section{General budgets by removing a mode}\label{sec:general-budgets}

Bounds for every block budget need no reach theorem beyond four blocks;
the exact reach results above serve only as seeds. The bounds rest on the
following mode-removal principle for arbitrary measurable inputs. It
preserves distinctness of the selected modes and applies even when the
one-sided coefficient is smaller than the reciprocal of the number of
modes.

\begin{lemma}[Mode removal]\label{lem:mode-removal}
Let $n\ge3$ and $k\ge2$. Suppose that, on every positive horizon, every
$(n-1)$-mode input admits at most $k-1$ distinct activation blocks with
one-sided error at most $C$ times its horizon, where $C>0$. Then the
corresponding $n$-mode, $k$-block coefficient is
\begin{equation}\label{eq:mode-removal}
 f_n(C)=\frac{(n-1)^2C+1}{n(n-1)(1+C)}.
\end{equation}
One may remove a mode of largest terminal mass if $C\ge1/(n-1)$, or
one of smallest terminal mass if $C\le1/(n-1)$. At equality either choice
is valid.
\end{lemma}
\begin{proof}
Set $E=f_n(C)T>0$, choose the indicated mode $q$, and write $a=A_q(T)$.
If $a\ge T-E$, use $q$ throughout. Its one-sided discrepancy is
$t-A_q(t)$, which is nondecreasing and ends at $T-a\le E$; every
other mode has nonpositive discrepancy $W_i-A_i$. This branch uses one
block. Otherwise put
\begin{equation}\label{eq:removal-length}
 L=T-a-E>0,\qquad E'=E-\frac{a}{n-1}.
\end{equation}
We verify both the positivity of $E'$ and the inequality $CL\le E'$.
A useful identity is
\begin{equation}\label{eq:removal-sign}
 f_n(C)-\frac1n
 =\frac{(n-2)((n-1)C-1)}{n(n-1)(1+C)}.
\end{equation}
If $C\ge1/(n-1)$, then $E\ge T/n$ and
$a<T-E\le(n-1)E$, so $E'>0$. Moreover $a\ge T/n$ and the function
\[
 a\longmapsto\frac{CT+(1/(n-1)-C)a}{1+C}
\]
is nonincreasing. Its value at $T/n$ is exactly $E$, whence
\begin{equation}\label{eq:removal-contract}
 E\ge\frac{CT+(1/(n-1)-C)a}{1+C}
 \quad\Longleftrightarrow\quad CL\le E'.
\end{equation}
If $0<C<1/(n-1)$, then $a\le T/n$ and
\[
 \frac{T}{n(n-1)}<E<\frac Tn.
\]
The first strict inequality follows by subtracting $1/[n(n-1)]$ from
$f_n(C)$; the numerator of the difference is $n(n-2)C>0$ before
cancelling the common denominator. Thus $E'>0$, and also
$L>T(1-2/n)>0$, so the constant-mode branch cannot occur. The same
function of $a$ is now increasing, and $a\le T/n$ again proves
\eqref{eq:removal-contract}. At $C=1/(n-1)$ the function is constant;
for minimum-mass selection $a\le T/n<(n-1)E$ also verifies positivity.

For $i\ne q$ complete the remaining rates on $[0,L]$ by
\[
 \widehat\alpha_i(t)=\alpha_i(t)+\frac{\alpha_q(t)}{n-1},
 \qquad \widehat A_i(t)=A_i(t)+\frac{A_q(t)}{n-1}.
\]
They are measurable, nonnegative, and sum to one. Apply the assumed
$(k-1)$-block construction to this input on $[0,L]$. Its error against
the original cumulative allocations satisfies, for every $t\le L$,
\[
 W_i(t)-A_i(t)
 =W_i(t)-\widehat A_i(t)+\frac{A_q(t)}{n-1}
 \le CL+\frac a{n-1}\le E.
\]
Mode $q$ has no service on this prefix and therefore has nonpositive
one-sided discrepancy. Append $q$ on $[L,T]$. Its discrepancy
$t-L-A_q(t)$ is nondecreasing and ends at $T-L-a=E$. Every other
one-sided discrepancy is nonincreasing after $L$. The prefix avoids
$q$, so appending it preserves distinctness. These arguments use
absolute continuity and monotonicity, and do not require piecewise
constant input rates.
\end{proof}

The fractional-linear recurrence has a simple coordinate. For a coefficient
$c$ in $j$ modes define
\[
 \eta_j=\frac{jc-1}{jc+1},\qquad
 c=\frac{1+\eta_j}{j(1-\eta_j)}.
\]
Positive $c$ is equivalent to $-1<\eta_j<1$. Substitution in
\eqref{eq:mode-removal} gives
\begin{equation}\label{eq:eta-transfer}
 \eta_j=\frac{j-2}{j}\eta_{j-1}.
\end{equation}
In particular a negative transformed seed remains negative. Replacing it
by zero is unnecessary and weakens the one-sided bound.

\begin{theorem}[Elementary coefficient for every spare-mode budget]
\label{thm:general-coefficient}
For $1\le k<n$, every $A\in\Acal_n(T)$ admits at most $k$ distinct
activation blocks with one-sided error at most $C_{n,k}T$, where
\begin{equation}\label{eq:elementary-coefficient}
 C_{n,k}=\frac{n(n-1)+(n-k)(n-k-1)}{nk(2n-k-1)}.
\end{equation}
In particular $C_{n,k}\ge1/n$, with equality exactly when $k=n-1$.
\end{theorem}
\begin{proof}
For $k=1$, choose a mode with mass at least $T/n$ and use it constantly.
Its error is at most $(n-1)T/n=C_{n,1}T$. For $k>1$ use
Lemma~\ref{lem:mode-removal} and induction at $(n-1,k-1)$. To solve
the recurrence put $m=n-k+1\ge2$. The single-block seed has
$mc_m=m-1$ and hence $\eta_m=(m-2)/m$. Iterating
\eqref{eq:eta-transfer}, including this seed factor, gives
\[
 \eta_n=\prod_{j=m}^n\frac{j-2}{j}
       =\frac{(n-k)(n-k-1)}{n(n-1)}.
\]
The identity also holds when $m=2$, since both sides are zero. Its value
lies in $[0,1)$; inversion gives~\eqref{eq:elementary-coefficient}.
The denominator is positive because $k\ge1$ and $n>k$. The assertion
about $1/n$ follows from the sign and zeros of $\eta_n$. This also
verifies that the maximum-mass branch alone suffices for this elementary
induction, though the more general lemma will be needed for sharper seeds.
\end{proof}

\begin{corollary}[One-sided, equal-mass, and unrestricted bounds]
\label{cor:general-bounds}
Write
\begin{equation}\label{eq:uniform-lower-general}
 L_{n,k}=\frac1{n(r_n^k-1)}.
\end{equation}
For $1\le k<n$ one has
\begin{align}
 TL_{n,k}&\le G^-_{n,k}(T)\le TC_{n,k},\label{eq:general-one-sided}\\
 T\max\{1/n,L_{n,k}\}&\le G^=_{n,k}(T)\le TC_{n,k},
       \label{eq:general-equal}\\
 T\max\{1/(k+1),L_{n,k}\}&\le F_{n,k-1}(T)
 \le T\max\{1/(k+1),C_{n,k}\},\label{eq:general-full}
\end{align}
where $G^=_{n,k}(T)$ denotes the worst full error over inputs with all
terminal masses $T/n$, using at most $k$ blocks. More generally, the
schedule in Theorem~\ref{thm:general-coefficient} has full error at most
$TC_{n,k}$ whenever every $m_i\le TC_{n,k}$.
\end{corollary}
\begin{proof}
For the uniform input, even a schedule with repetitions satisfies the
block-end recurrence $t_j\le r_n(t_{j-1}+E)$ when its one-sided error
is $E$. Iterating gives $T\le nE(r_n^k-1)$. Distinct greedy blocks attain
this lower coefficient; their full uniform error is
$T\max\{1/n,L_{n,k}\}$ by Theorem~\ref{thm:uniform}.
This input belongs to all three input classes. Every positive discrepancy
satisfies $A_i(t)-W_i(t)\le m_i$, proving the mass-qualified and equal-mass
upper bounds. Finally Theorem~\ref{thm:one-sided-reduction} proves both
unrestricted bounds from the one-sided ones and the omitted-mode threshold.
\end{proof}

\begin{corollary}[An exact plateau]\label{cor:general-plateau}
For integers $k\ge2$ and $k+1\le n\le k(k+1)/2$,
\begin{equation}\label{eq:general-plateau}
 F_{n,k-1}(T)=\frac{T}{k+1}.
\end{equation}
This is a guaranteed plateau interval, not a claim that it is maximal.
\end{corollary}
\begin{proof}
Direct subtraction and factorization give
\begin{equation}\label{eq:plateau-factor}
 C_{n,k}-\frac1{k+1}
 =\frac{2(n-k-1)(n-k(k+1)/2)}{nk(k+1)(2n-k-1)}.
\end{equation}
The denominator is positive. The numerator is nonpositive on the stated
interval, so~\eqref{eq:general-full} proves equality. Its lower witness
is $k+1$ successive distinct pure modes of duration $T/(k+1)$.
For $k=1$ the interval in the statement is empty; separately
$C_{n,1}=(n-1)/n$, with equality to $1/2$ at $n=2$.
\end{proof}

\begin{proposition}[The first mode-count correction]\label{prop:general-asymptotics}
For fixed $k\ge1$ and $n\to\infty$,
\begin{align}
 C_{n,k}&=\frac1k-\frac{k+1}{2kn}
                 +\frac{k^2-1}{4kn^2}+O_k(n^{-3}),
                 \label{eq:elementary-series}\\
 L_{n,k}&=\frac1k-\frac{k+1}{2kn}
                 +\frac{k^2-1}{12kn^2}+O_k(n^{-3}).
                 \label{eq:lower-series}
\end{align}
Each of $G^-_{n,k}(T)$, $G^=_{n,k}(T)$, and $F_{n,k-1}(T)$ consequently
equals
\begin{equation}\label{eq:sharp-first-correction}
 \frac Tk-\frac{(k+1)T}{2kn}+O_k(Tn^{-2}).
\end{equation}
The gap between the displayed elementary upper and uniform lower
coefficients is $(k^2-1)/(6kn^2)+O_k(n^{-3})$.
\end{proposition}
\begin{proof}
Put $x=1/n$. Dividing numerator and denominator of
\eqref{eq:elementary-coefficient} by $n^2$ expresses it as
\[
 \frac{2-2(k+1)x+k(k+1)x^2}{k(2-(k+1)x)}.
\]
Expand the denominator using the geometric series through degree two.
For the lower coefficient use
\[
 (1-x)^{-k}=1+kx+\frac{k(k+1)}2x^2
             +\frac{k(k+1)(k+2)}6x^3+O_k(x^4)
\]
and invert $[(1-x)^{-k}-1]/x$. These calculations give
\eqref{eq:elementary-series}--\eqref{eq:lower-series}, with bounded
remainders near $x=0$ because the resulting rational denominators do not
vanish there. For large $n$, both coefficients tend to $1/k>1/(k+1)$
and $1/n$ is smaller. The three pairs of bounds in
Corollary~\ref{cor:general-bounds} therefore give the stated expansions.
When $k=1$ both coefficient formulas equal $1-1/n$ exactly; their
second-order terms and gap vanish.
\end{proof}

\subsection{Propagating exact seeds}\label{subsec:seeded}

\begin{theorem}[A general seed and the four-block improvement]
\label{thm:seeded}
Let $m\ge2$ and $\ell\ge1$ be integers. Suppose an $m$-mode one-sided coefficient $c_m>0$ is proved using at most
$\ell$ distinct blocks. For $n\ge m$, $k=\ell+n-m$, the propagated
coefficient is
\begin{equation}\label{eq:general-seed}
 C^{[c_m]}_{n,k}=\frac{1+\theta}{n(1-\theta)},\qquad
 \theta=\frac{m(m-1)}{n(n-1)}\frac{mc_m-1}{mc_m+1}.
\end{equation}
In particular, for $1\le\ell\le4$ and $n>k\ge\ell$, set $m=n-k+\ell$
and use the established seed $c_m=L_{m,\ell}$. This yields
\begin{equation}\label{eq:exact-seed-theta}
 C^{(\ell)}_{n,k}=\frac{1+\theta^{(\ell)}_{n,k}}
                         {n(1-\theta^{(\ell)}_{n,k})},\qquad
 \theta^{(\ell)}_{n,k}=\frac{m(m-1)}{n(n-1)}
             \left[2\left(\frac{m-1}{m}\right)^\ell-1\right].
\end{equation}
Write $C^*_{n,k}=C^{(4)}_{n,k}$ for $k\ge4$. Then
\begin{equation}\label{eq:seeded-full}
 G^-_{n,k}(T)\le TC^*_{n,k},\qquad
 F_{n,k-1}(T)\le T\max\{1/(k+1),C^*_{n,k}\}.
\end{equation}
For every $n>k\ge4$, $C^*_{n,k}<C_{n,k}$. The four-block seed, and thus
this improvement, retains the certificate dependency of
Theorem~\ref{thm:reach-four}.
\end{theorem}
\begin{proof}
Apply Lemma~\ref{lem:mode-removal} $n-m$ times. Equation
\eqref{eq:eta-transfer} telescopes to the factor $m(m-1)/[n(n-1)]$.
The seed transform lies strictly between $-1$ and one, and multiplication
by a factor in $(0,1]$ preserves this interval. Thus all coefficients are
positive and all denominators are positive. Distinctness is preserved at
every step. Substituting $c_m=L_{m,\ell}$ gives the bracket in
\eqref{eq:exact-seed-theta}. The seed is available for $m\ge\ell+1$
by the established reach theorems. The heavy-mode reduction proves the
full bound.

For strict comparison, the elementary construction at $k=4$ satisfies
\[
 C_{m,4}-L_{m,4}
 =\frac{10m^2-5m+1}
 {2(2m-5)(2m-1)(2m^2-2m+1)}>0\qquad(m\ge5).
\]
All denominator factors are positive, and the numerator is
$5m(2m-1)+1>0$. Further,
$f_j'(c)=(j-2)/[(j-1)(1+c)^2]>0$ for $j\ge3$.
The elementary coefficient is the same recurrence from its elementary
four-block seed; strict improvement therefore persists at every step.
This comparison is about one-sided coefficients: the maxima in the full
bounds can coincide on a plateau.
\end{proof}

\begin{corollary}[Closed consequences of the seed]\label{cor:seed-consequences}
For $k\ge4$, equal-terminal-mass inputs have exact worst full error $T/n$
when $n=k+1$ or $n=k+2$. In addition,
\begin{equation}\label{eq:extra-plateau}
 F_{16,4}(T)=T/6.
\end{equation}
For all $n>k\ge4$, a sufficient exact-plateau criterion, involving only
integers, is
\begin{equation}\label{eq:seed-plateau-test}
 (n+k+1)(m-1)[2(m-1)^4-m^4]
 \le(n-k-1)n(n-1)m^3,\qquad m=n-k+4.
\end{equation}
\end{corollary}
\begin{proof}
The fourth-power seed transform is negative exactly at the integer mode
counts $m=5,6$: $2(4/5)^4<1$, $2(5/6)^4<1$, and
$2(6/7)^4>1$, while $(1-1/m)^4$ increases with $m$. Hence $C^*_{n,k}<1/n$
on the two specified diagonals. The seeded schedule then has full error
at most $T/n$ for equal terminal masses. Every schedule with $k<n$ blocks
omits a mode, so the uniform input gives the matching $T/n$ lower bound.
The broader certificate-free equal-mass band in
Corollary~\ref{cor:light-boundaries} also includes both diagonals. Retaining
the negative transform supplies the stronger one-sided coefficient below
$1/n$, even though clipping would suffice for these full-error values.

At $(n,k)=(16,5)$ direct evaluation gives
\[
 C^*_{16,5}=\frac{588449}{3544816}<\frac16
       <\frac{35}{208}=C_{16,5}.
\]
Six successive pure modes give the lower bound $T/6$. This extends the
elementary four-switch plateau through one additional mode count.
At $n=17,k=5$ the seeded coefficient exceeds $1/6$, so this bound
alone certifies no further extension there.
Finally $C^*_{n,k}\le1/(k+1)$ is equivalent to
$(n+k+1)\theta^{(4)}_{n,k}\le n-k-1$. Substituting the expression for
$\theta$ and clearing the positive denominator $n(n-1)m^3$ gives
\eqref{eq:seed-plateau-test}. Whenever it holds, the full upper bound
matches the omitted-mode lower witness. It characterizes this certificate's
plateau criterion, not every parameter at which the actual minimax may
have that value.
\end{proof}

\begin{proposition}[Seed-dependent second-order gap]\label{prop:seed-series}
For fixed $k\ge\ell$, $1\le\ell\le4$, as $n\to\infty$,
\begin{align}
 C^{(\ell)}_{n,k}
 &=\frac1k-\frac{k+1}{2kn}
 +\frac{3k^3-3k-2\ell^3+2\ell}{12k^2n^2}+O_{k,\ell}(n^{-3}),
 \label{eq:seed-series}\\
 C^{(\ell)}_{n,k}-L_{n,k}
 &=\frac{(k-\ell)(k^2+k\ell+\ell^2-1)}{6k^2n^2}
       +O_{k,\ell}(n^{-3}).\label{eq:seed-gap}
\end{align}
For $\ell=4$ the latter coefficient is $(k^3-k-60)/(6k^2)$.
For $k=5$ it is $2/5$, half the elementary gap coefficient $4/5$.
These are bounds on the unresolved second-order error, not an exact
second-order expansion of the general minimax.
\end{proposition}
\begin{proof}
With $x=1/n$ and $d=k-\ell$, the transformed coefficient is
\[
 \frac{(1-dx)(1-(d+1)x)}{1-x}
 \left[2\left(\frac{1-(d+1)x}{1-dx}\right)^\ell-1\right].
\]
Expanding through degree three gives
\[
 \theta=1-2kx+k(k-1)x^2+
 \left(k(k-1)-\frac{\ell^3-\ell}{3}\right)x^3+O_{k,\ell}(x^4).
\]
Substitute into $x(1+\theta)/(1-\theta)$ and cancel $x$ from the
denominator. Its new constant term is $2k>0$, so ordinary power-series
division is valid and gives~\eqref{eq:seed-series}. Subtracting
\eqref{eq:lower-series} and using
$k^3-k-\ell^3+\ell=(k-\ell)(k^2+k\ell+\ell^2-1)$ gives the gap.
For $k=\ell$ the seed is exactly $L_{n,k}$, so in that case the entire
gap is zero, not merely its displayed second-order coefficient.
\end{proof}
